\begingroup
\appendix
\ctexset{chapter={name={附录,},number=\Alph{chapter},
  tocline={\protect\numberline{\thechapter}#2}}}
\hypersetup{bookmarksnumbered=true}

% 只调整目录及书签：附录作为父项，A--F作为次一级条目。
% 正文仍使用chapter及原有编号；各附录内部的节不写入目录。
\makeatletter
\let\mathnotes@appendix@addcontentsline\addcontentsline
\def\mathnotes@appendix@first{first}
\def\mathnotes@appendix@toc{toc}
\def\mathnotes@appendix@chapter{chapter}
\renewcommand\addcontentsline[3]{%
  \def\mathnotes@appendix@file{#1}%
  \def\mathnotes@appendix@level{#2}%
  \ifx\mathnotes@appendix@file\mathnotes@appendix@toc
    \ifx\mathnotes@appendix@level\mathnotes@appendix@chapter
      \ifx\mathnotes@appendix@first\@empty\else
        \edef\mathnotes@appendix@href{\@currentHref}%
        \phantomsection
        \mathnotes@appendix@addcontentsline{toc}{chapter}{附录}%
        \let\@currentHref\mathnotes@appendix@href
        \let\mathnotes@appendix@first\@empty
      \fi
      \mathnotes@appendix@addcontentsline{toc}{section}{#3}%
    \fi
  \else
    \mathnotes@appendix@addcontentsline{#1}{#2}{#3}%
  \fi}
\makeatother

\chapter{集合与常用逻辑用语}\label{app:sets-logic}
\begingroup
\fontsize{12pt}{20pt}\selectfont
\setlength{\parskip}{4pt}

\section{集合}
把所讨论的对象看成一个整体，称为{\heiti{}集合}，其中的每一个对象称为集合的{\heiti{}元素}。\(x\in A\) 表示 \(x\) 属于 \(A\)，\(x\notin A\) 表示它不属于 \(A\)。给出一个集合，要能确定哪些对象属于它。例如，“小于 \(5\) 的正整数”确定一个集合，“很小的正整数”则还需说明多小算小。

将元素逐个写出叫作{\heiti{}列举法}，如 \(\{1,2,3,4\}\)；规定元素的范围和条件叫作{\heiti{}描述法}，同一集合可以写成 \(\{n\in\mathbb N^*:n<5\}\)，冒号读作“使得”。本书用 \(\mathbb N,\mathbb N^*,\mathbb Z,\mathbb Q,\mathbb R\) 分别表示非负整数、正整数、整数、有理数、实数组成的集合。

\begin{definition}
若 \(A\) 的每个元素都是 \(B\) 的元素，称 \(A\) 是 \(B\) 的{\heiti{}子集}，记为 \(A\subseteq B\)。若两集合的元素完全相同，称它们相等，记为 \(A=B\)；若 \(A\subseteq B\) 且 \(A\ne B\)，则称 \(A\) 是 \(B\) 的{\heiti{}真子集}，记为 \(A\subsetneq B\)。
\end{definition}
所以证明 \(A=B\)，可以分别证明 \(A\subseteq B\) 和 \(B\subseteq A\)。元素的次序和书写次数均不改变集合，如 \(\{1,2\}=\{2,1,2\}\)。没有元素的集合叫作{\heiti{}空集}，记为 \(\varnothing\)。它是每个集合的子集，因为其中没有一个元素能够违反包含关系。注意 \(1\in\{1,2\}\) 说的是元素的归属，\(\{1\}\subseteq\{1,2\}\) 说的是两个集合的包含关系。

\begin{example}
写出所有既是整数又满足 \(x^2<5\) 的数。若把“整数”改为“实数”，所得集合还相同吗？
\end{example}

\begin{exercise}
\mbox{}\par\nobreak

（1）设 \(A=\{x\in\mathbb Z:x=2k+1,\ k\in\mathbb Z\}\)，\(B=\{x\in\mathbb Z:x=2k-1,\ k\in\mathbb Z\}\)。证明 \(A=B\)。

（2）列出 \(\{0,1\}\) 的全部子集。其中哪些以 \(0\) 为元素？
\end{exercise}

\section{集合的运算}
设 \(A,B\) 是两个集合。把同时属于两者的元素取出，得到{\heiti{}交集} \(A\cap B\)；把至少属于其一的元素取出，得到{\heiti{}并集} \(A\cup B\)；从 \(A\) 中去掉属于 \(B\) 的元素，得到{\heiti{}差集} \(A\setminus B\)。交与并交换两个集合不影响结果，差集则要分清从哪个集合中取元素。

在数轴上，一段连续取值的范围常用区间表示。以 \(a<b\) 为例，\([a,b]=\{x\in\mathbb R:a\le x\le b\}\)，\((a,b)=\{x\in\mathbb R:a<x<b\}\)。方括号表示取到相应端点，圆括号表示不取；例如 \((-\infty,b]\) 表示所有不大于 \(b\) 的实数。

事先确定讨论的全体对象组成的集合 \(\Omega\)，称为{\heiti{}全集}。对 \(A\subseteq\Omega\)，集合 \(\Omega\setminus A\) 称为 \(A\) 在 \(\Omega\) 中的{\heiti{}补集}。同是偶数集，在整数集中的补集是奇数集，在实数集中的补集还含有 \(\tfrac12\) 等非整数。
\begin{proposition}
设 \(A,B\subseteq\Omega\)，则 \(\Omega\setminus(A\cup B)=(\Omega\setminus A)\cap(\Omega\setminus B)\)。
\end{proposition}
\begin{proof}
任取 \(x\in\Omega\)。它属于左边，恰好是指它既不属于 \(A\)，也不属于 \(B\)；这又恰好是指它同时属于两个补集。因此等式两边有完全相同的元素。
\end{proof}

还有一种不同的运算。分别取 \(a\in A,b\in B\)，把两者依次配成有序对 \((a,b)\)。所有这样的有序对组成的集合称为{\heiti{}笛卡尔积} \(A\times B\)。平面直角坐标系中的坐标对全体便是 \(\mathbb R\times\mathbb R\)；交换两个坐标一般会得到另一个点。

\begin{example}
设 \(A=(-2,1]\)，\(B=[1,3)\)。求 \(A\cap B,\ A\cup B,\ A\setminus B\)，并说明 \(1\) 在三个集合中分别是否出现。
\end{example}

\begin{exercise}
\mbox{}\par\nobreak

（1）设 \(A,B\subseteq\Omega\)。证明 \(\Omega\setminus(A\cap B)=(\Omega\setminus A)\cup(\Omega\setminus B)\)。

（2）设 \(A=\{0,1\}\)，\(B=\{1,2\}\)。列出 \(A\times B\) 的元素，并求 \((A\times B)\cap(B\times A)\)。
\end{exercise}

\section{命题与条件}
可以判断真假的陈述句称为{\heiti{}命题}。例如“\(6\) 是偶数”是真命题，“所有的奇数都是质数”是假命题。“\(n\) 是偶数”还要看 \(n\) 取什么值；它给出了一个关于 \(n\) 的条件，可记为 \(P(n)\)。

对命题 \(P,Q\)，\(\neg P\) 读作“非 \(P\)”，其真假与 \(P\) 相反；\(P\land Q\) 要求二者都真，\(P\lor Q\) 要求至少一个为真，也允许二者都真。因此否定“二者都真”，只须说至少一个不真；否定“至少一个真”，则须说二者都不真。这正是上一节两条补集公式中“且”与“或”互换的原因。

\begin{definition}
若 \(P\) 成立时 \(Q\) 必成立，记为 \(P\Rightarrow Q\)，称 \(P\) 是 \(Q\) 的{\heiti{}充分条件}，\(Q\) 是 \(P\) 的{\heiti{}必要条件}。若两方向均成立，记为 \(P\Leftrightarrow Q\)，称二者互为{\heiti{}充要条件}。
\end{definition}
给定范围 \(\Omega\)，设满足条件 \(P,Q\) 的元素分别组成 \(A,B\)。那么“对每个 \(x\in\Omega\)，由 \(P(x)\) 能推出 \(Q(x)\)”恰好是 \(A\subseteq B\)。若要推翻它，只需指出一个属于 \(A\) 却不属于 \(B\) 的元素。这就是{\heiti{}反例}：它必须满足前提，同时违反结论。

交换前提和结论得到{\heiti{}逆命题} \(Q\Rightarrow P\)，未必与原命题同真同假。所有 \(4\) 的倍数都是偶数，但偶数 \(2\) 不是 \(4\) 的倍数。若同时否定前提和结论再交换位置，得到{\heiti{}逆否命题} \(\neg Q\Rightarrow\neg P\)。它与原命题等价，因为让两者不成立所需的反例完全相同：\(P\) 真而 \(Q\) 假。

\begin{example}
设 \(x\) 为实数。“\(x>1\)”是“\(x^2>1\)”的什么条件？增加对 \(x\) 的范围限制，使两者等价。
\end{example}

\begin{exercise}
\mbox{}\par\nobreak

（1）设 \(a,b\) 是整数。写出“若 \(a+b\) 是奇数，则 \(a,b\) 中至少一个是奇数”的逆命题与逆否命题，并判断真假。

（2）对于两个集合 \(A,B\)，证明 \(A\subseteq B\) 的一个充要条件是 \(A\cap B=A\)。
\end{exercise}

\section{量词}
有两种常用的方法，可以从条件 \(P(x)\) 得到一个命题。\(\forall x\in A,\ P(x)\) 读作“对 \(A\) 中每个 \(x\)，\(P(x)\) 都成立”，其中 \(\forall\) 称为{\heiti{}全称量词}；\(\exists x\in A,\ P(x)\) 读作“在 \(A\) 中存在 \(x\)，使 \(P(x)\) 成立”，其中 \(\exists\) 称为{\heiti{}存在量词}。这里的“存在”指至少有一个；若还要求只有一个，须另证唯一性。

量词涉及哪一范围，也是命题的一部分。例如，“存在 \(x\) 使 \(x^2=2\)”在整数集中为假，在实数集中为真。几个量词连用时，要从左向右读：后选的量可以依赖先给的量，先选的量必须对后面的选择都有效。

\begin{example}
比较“每个整数都有比它大的整数”和“有一个整数比所有整数都大”。用量词写出两句话，并分别判断真假。
\end{example}

要否定一个全称命题，只须找出一个例外；要否定一个存在命题，则要排除全部可能。因此
\[
 \neg(\forall x\in A,\ P(x))\ \Leftrightarrow\ \exists x\in A,\ \neg P(x),\qquad
 \neg(\exists x\in A,\ P(x))\ \Leftrightarrow\ \forall x\in A,\ \neg P(x).
\]
有多层量词时逐层否定，原有的先后次序不变。例如，“每本书都有人读过”的否定是“有一本书谁都没读过”。

学到第一章的式\eqref{eq:analysis-sequence-limit-definition}时，可以再用一次这个规则。数列收敛于 \(a\)，要求对每个 \(\varepsilon>0\) 都能选定 \(N\)，使所有 \(n>N\) 都满足 \(|a_n-a|<\varepsilon\)。它的否定是：存在一个固定的 \(\varepsilon_0>0\)，无论给出多大的 \(N\)，总能找到某个 \(n>N\)，使 \(|a_n-a|\geq\varepsilon_0\)。改变的是量词和末尾的不等式，\(N\) 与 \(n\) 的选择顺序始终相同。

\begin{exercise}
\mbox{}\par\nobreak

（1）写出“对每个实数 \(x\)，都有 \(x^2\geq x\)”的否定，并给出一个使否定成立的数。

（2）已知整数集 \(S\) 非空。用量词写出“\(S\) 有最大元素”，再把它的否定译成一句话。
\end{exercise}
\endgroup

\chapter{常用记号}
\label{app:notation}

本附录汇集正文常用的记号；各题中的定义域和约定以题设为准。

\section{数集与运算记号}
\(\mathbb N,\mathbb N^*,\mathbb Z,\mathbb Q,\mathbb R,\mathbb C\) 分别表示非负整数集、正整数集、整数集、有理数集、实数集和复数集。有限集 \(A\) 的元素个数记为 \(|A|\)。集合运算及逻辑记号见附录\ref{app:sets-logic}。

\(\sum_{k=1}^{n}a_k=a_1+\cdots+a_n\)，\(\prod_{k=1}^{n}a_k=a_1\cdots a_n\)。两式中 \(k\) 均依次取 \(1\) 至 \(n\) 的整数。

\section{函数与线性运算}

\(f:A\to B\)表示映射的定义域为\(A\)、陪域为\(B\)；其像为\(f(A)=\{f(x):x\in A\}\)。\((g\circ f)(x)=g(f(x))\)表示先使用\(f\)，再使用\(g\)。\(f^{-1}\)表示已经存在的反函数，与逐点取倒数\(1/f\)分别理解。

线性映射\(T:V\to W\)的核与像分别为
\begin{equation*}
 \ker T=\{v\in V:Tv=0\},\qquad
 \operatorname{im}T=\{Tv:v\in V\}.
\end{equation*}
\(\operatorname{span}(v_1,\ldots,v_k)\)表示这些向量全部线性组合构成的子空间。内积记为\(\langle u,v\rangle\)，其诱导的长度记为\(\|u\|=\sqrt{\langle u,u\rangle}\)；几何向量的长度在正文中也记为\(|u|\)。

矩阵的转置为\(A^T\)，可逆矩阵的逆为\(A^{-1}\)，单位矩阵为\(I\)。行列式记为\(\det A\)。这些记号的作用对象由当前的向量空间与坐标表示确定。

\section{极限与近似}

\(a_n\to a\)表示数列收敛到\(a\)，\(f(x)\to L\)、\(x\to x_0\)表示相应函数极限。余项\(r(h)=o(h^n)\)表示\(r(h)/h^n\to0\)；\(r(h)=O(h^n)\)表示在所讨论极限附近，\(|r(h)|\le C|h|^n\)可由某个固定常数\(C\)控制。

\(C[a,b]\)、\(C^1[a,b]\)、\(C^2[a,b]\)依次表示连续函数、一阶连续可导函数、二阶连续可导函数组成的空间。连续函数的最大范数为\(\|f\|_\infty=\max_{[a,b]}|f|\)，比较两个函数时由\(\|f-g\|_\infty\)衡量整个区间上的最大误差。

\chapter{常用不等式}
\label{app:inequalities}

本附录整理经常用于估计、比较和误差控制的不等式，并列出其主要依据。选择不等式时，应注意当前需要控制的量以及等号对应的结构。

\section{Cauchy不等式}

对实数\(a_1,\ldots,a_n\)、\(b_1,\ldots,b_n\)，
\begin{equation*}
 \left(\sum_{i=1}^n a_ib_i\right)^2
 \le\left(\sum_{i=1}^n a_i^2\right)\left(\sum_{i=1}^n b_i^2\right).
\end{equation*}
等号当且仅当两个向量\((a_i)\)、\((b_i)\)线性相关。可由恒等式
\begin{equation*}
 \left(\sum a_i^2\right)\left(\sum b_i^2\right)-\left(\sum a_ib_i\right)^2
 =\sum_{i<j}(a_ib_j-a_jb_i)^2
\end{equation*}
直接得到。内积形式为\(|\langle u,v\rangle|\le\|u\|\|v\|\)，从而
\begin{equation*}
 \|u+v\|\le\|u\|+\|v\|,\qquad
 \bigl|\|u\|-\|v\|\bigr|\le\|u-v\|.
\end{equation*}

对\(x_i>0\)，取\(a_i=\sqrt{x_i}\)、\(b_i=1/\sqrt{x_i}\)，可得
\begin{equation*}
 \left(\sum_{i=1}^n x_i\right)\left(\sum_{i=1}^n\frac1{x_i}\right)\ge n^2.
\end{equation*}

\section{平均值不等式}

对正数\(a_1,\ldots,a_n\)，调和平均、几何平均、算术平均与平方平均满足
\begin{equation*}
 \frac{n}{\sum_{i=1}^n1/a_i}
 \le\left(\prod_{i=1}^n a_i\right)^{1/n}
 \le\frac1n\sum_{i=1}^n a_i
 \le\sqrt{\frac1n\sum_{i=1}^n a_i^2}.
\end{equation*}
各处等号均当且仅当全部\(a_i\)相等。

加权形式中，设\(w_i>0\)、\(\sum w_i=1\)，则
\begin{equation*}
 \prod_{i=1}^n a_i^{w_i}\le\sum_{i=1}^n w_i a_i.
\end{equation*}
取\(A=\sum w_i a_i\)，对每个\(a_i/A\)使用\(\ln x\le x-1\)，加权相加便得\(\sum w_i\ln(a_i/A)\le0\)，即上式。对\(1/a_i\)应用同一结论得到前面的调和平均界，平方平均界则由Cauchy不等式得到。

\section{指数、对数与幂}

常用的切线估计为
\begin{equation*}
 e^x\ge1+x\quad(x\in\mathbb R),\qquad
 \ln x\le x-1\quad(x>0).
\end{equation*}
等号分别在\(x=0\)、\(x=1\)处取得。由对数的积分表达或作差求导，\(x>0\)时还有
\begin{equation*}
 \frac{x}{1+x}<\ln(1+x)<x.
\end{equation*}

Bernoulli不等式为
\begin{equation*}
 (1+x)^n\ge1+nx\qquad(x\ge-1,\ n\in\mathbb N^*),
\end{equation*}
可以直接归纳证明。对于实数指数\(r\)与\(x>-1\)，由函数的凸性可得
\begin{equation*}
 \begin{cases}
 (1+x)^r\ge1+rx,&r\le0\text{或}r\ge1,\\
 (1+x)^r\le1+rx,&0\le r\le1.
 \end{cases}
\end{equation*}

\section{函数增量与积分估计}

若\(f\)在连接\(x,y\)的闭区间连续、内部可导，且\(|f'|\le M\)，则中值定理给出
\begin{equation*}
 |f(x)-f(y)|\le M|x-y|.
\end{equation*}
若\(f\)在\([a,b]\)可积，\(m\le f\le M\)，则
\begin{equation*}
 m(b-a)\le\int_a^bf(x)\,dx\le M(b-a),\qquad
 \left|\int_a^bf\right|\le\int_a^b|f|.
\end{equation*}
特别地，连续函数满足
\begin{equation*}
 \left|\int_a^b(f-g)\right|\le(b-a)\|f-g\|_\infty.
\end{equation*}
这些估计分别把局部变化率和整个区间上的函数误差，转化为可以计算的数值界。

\chapter{常用公式}
\label{app:formulas}

\begingroup
\fontsize{12pt}{20pt}\selectfont
\setlength{\parskip}{4pt}
\setlength{\abovedisplayskip}{7pt}
\setlength{\belowdisplayskip}{7pt}
\setlength{\abovedisplayshortskip}{4pt}
\setlength{\belowdisplayshortskip}{4pt}
\setlength{\LTpre}{6pt}
\setlength{\LTpost}{6pt}
\setlength{\jot}{1pt}

以下公式按对象集中列出，具体结论的定义、证明与使用见正文。求导、积分及三角函数公式中的角度均采用弧度制。

\section{数列与极限}
\begin{equation*}
 \begin{aligned}
 &a_n=a_1+(n-1)d,&
 &S_n=\frac{n(a_1+a_n)}2,\\
 &a_n=a_1q^{n-1},&
 &S_n=\frac{a_1(1-q^n)}{1-q}\quad(q\ne1).
 \end{aligned}
\end{equation*}
等比数列\(q=1\)时\(S_n=na_1\)；\(|q|<1\)时\(\sum_{n=1}^\infty a_1q^{n-1}=a_1/(1-q)\)。由部分和恢复通项的关系为\(a_n=S_n-S_{n-1}\)，其中\(S_0=0\)。
\begin{equation*}
 \begin{aligned}
 \sum_{k=1}^n k&=\frac{n(n+1)}2,&
 \sum_{k=1}^n k^2&=\frac{n(n+1)(2n+1)}6,\\
 \sum_{k=1}^n k^3&=\left(\frac{n(n+1)}2\right)^2,&
 e&=\lim_{n\to\infty}\left(1+\frac1n\right)^n.
 \end{aligned}
\end{equation*}
基本函数极限有\(\lim_{x\to0}\sin x/x=1\)、\(\lim_{x\to0}(e^x-1)/x=1\)、\(\lim_{x\to0}\ln(1+x)/x=1\)。

\section{三角恒等式}
\begin{equation*}
 \begin{aligned}
 \sin(\alpha\pm\beta)&=\sin\alpha\cos\beta\pm\cos\alpha\sin\beta,\\
 \cos(\alpha\pm\beta)&=\cos\alpha\cos\beta\mp\sin\alpha\sin\beta,\\
 \tan(\alpha\pm\beta)&=\frac{\tan\alpha\pm\tan\beta}{1\mp\tan\alpha\tan\beta},\\
 \sin2\alpha&=2\sin\alpha\cos\alpha,\\
 \cos2\alpha&=\cos^2\alpha-\sin^2\alpha=1-2\sin^2\alpha=2\cos^2\alpha-1,\\
 \sin^2\alpha&=\frac{1-\cos2\alpha}2,\qquad
 \cos^2\alpha=\frac{1+\cos2\alpha}2.
 \end{aligned}
\end{equation*}
记\(m=(\alpha+\beta)/2\)、\(d=(\alpha-\beta)/2\)，则
\begin{equation*}
 \begin{aligned}
 \sin\alpha+\sin\beta&=2\sin m\cos d,&
 \sin\alpha-\sin\beta&=2\cos m\sin d,\\
 \cos\alpha+\cos\beta&=2\cos m\cos d,&
 \cos\alpha-\cos\beta&=-2\sin m\sin d.
 \end{aligned}
\end{equation*}
设\(R=\sqrt{a^2+b^2}>0\)，取\(\cos\varphi=a/R\)、\(\sin\varphi=b/R\)，则\(a\sin x+b\cos x=R\sin(x+\varphi)\)。三角形中，若边\(a,b,c\)分别对角\(A,B,C\)，外接圆半径为\(R_\triangle\)，则
\begin{equation*}
 \frac a{\sin A}=\frac b{\sin B}=\frac c{\sin C}=2R_\triangle,\qquad
 c^2=a^2+b^2-2ab\cos C,\qquad S=\frac12ab\sin C.
\end{equation*}

\section{向量与矩阵}

在正交单位基下，\(u\cdot v=\sum_i u_iv_i\)，\(\|u\|=\sqrt{\sum_i u_i^2}\)。非零向量满足
\begin{equation*}
 \cos\theta=\frac{u\cdot v}{\|u\|\|v\|},\qquad
 \operatorname{proj}_v u=\frac{u\cdot v}{v\cdot v}v.
\end{equation*}
在三维右手正交单位基下，叉积的坐标为
\begin{equation*}
 u\times v=(u_2v_3-u_3v_2,\ u_3v_1-u_1v_3,\ u_1v_2-u_2v_1).
\end{equation*}
平行四边形面积为\(\|u\times v\|\)，平行六面体体积为\(|u\cdot(v\times w)|\)。

\begin{equation*}
 A=\begin{pmatrix}a&b\\c&d\end{pmatrix},\qquad \det A=ad-bc,\qquad
 A^{-1}=\frac1{ad-bc}\begin{pmatrix}d&-b\\-c&a\end{pmatrix}\quad(ad-bc\ne0).
\end{equation*}
平面仿射变换\(x=Au+t\)的面积倍率为\(|\det A|\)。

\section{空间几何体}

记\(B,P\)为底面积、底面周长，\(h\)为高；\(B_1,B_2\)与\(P_1,P_2\)分别对应两底面。\(l\)表示正棱锥、正棱台的斜高，或圆锥、圆台的母线长。表面积计入底面。
\begingroup
\renewcommand{\arraystretch}{1.15}
\begin{longtable}{@{}p{.22\linewidth}p{.31\linewidth}p{.41\linewidth}@{}}
\toprule
几何体 & 体积\(V\) & 表面积\(S\) \\
\midrule
\endhead
正方体（棱长\(a\)） & \(a^3\) & \(6a^2\) \\
长方体（\(a,b,c\)） & \(abc\) & \(2(ab+bc+ca)\) \\
正四面体（棱长\(a\)） & \(\sqrt2\,a^3/12\) & \(\sqrt3\,a^2\) \\
直棱柱 & \(Bh\) & \(2B+Ph\) \\
正棱锥 & \(Bh/3\) & \(B+Pl/2\) \\
正棱台 & \(\dfrac h3(B_1+B_2+\sqrt{B_1B_2})\) & \(B_1+B_2+\dfrac{P_1+P_2}{2}l\) \\
圆柱（半径\(r\)） & \(\pi r^2h\) & \(2\pi r(r+h)\) \\
圆锥（半径\(r\)） & \(\pi r^2h/3\) & \(\pi r(r+l)\) \\
圆台（半径\(R,r\)） & \(\dfrac{\pi h}{3}(R^2+Rr+r^2)\) & \(\pi(R^2+r^2)+\pi(R+r)l\) \\
球（半径\(r\)） & \(4\pi r^3/3\) & \(4\pi r^2\) \\
半球（半径\(r\)） & \(2\pi r^3/3\) & \(3\pi r^2\)（含圆形底面） \\
\bottomrule
\end{longtable}
\endgroup
圆锥的\(l=\sqrt{h^2+r^2}\)，圆台的\(l=\sqrt{h^2+(R-r)^2}\)。上述棱柱、棱锥、棱台的体积公式也适用于一般情形。

\section{复数}
\begin{equation*}
 \overline{a+bi}=a-bi,\qquad
 |z|^2=z\overline z,\qquad
 \frac1z=\frac{\overline z}{|z|^2}\quad(z\ne0).
\end{equation*}
对非零复数\(z=r(\cos\theta+i\sin\theta)=re^{i\theta}\)，
\begin{equation*}
 z^n=r^n(\cos n\theta+i\sin n\theta),\qquad n\in\mathbb Z.
\end{equation*}
若\(w=\rho e^{i\varphi}\ne0\)、\(n\in\mathbb N^*\)，方程\(z^n=w\)的全部根为
\begin{equation*}
 z_k=\rho^{1/n}e^{i(\varphi+2k\pi)/n},\qquad k=0,1,\ldots,n-1.
\end{equation*}
Euler公式为\(e^{i\theta}=\cos\theta+i\sin\theta\)。

\section{二次曲线}

以下\(a,b,p>0\)，切线公式中的\((x_0,y_0)\)在相应曲线上。
\begingroup
\renewcommand{\arraystretch}{1.3}
\begin{longtable}{@{}p{.41\linewidth}p{.53\linewidth}@{}}
\toprule
标准方程 & 切线方程 \\
\midrule
\endhead
\(\dfrac{x^2}{a^2}+\dfrac{y^2}{b^2}=1\) & \(\dfrac{xx_0}{a^2}+\dfrac{yy_0}{b^2}=1\) \\[6pt]
\(\dfrac{x^2}{a^2}-\dfrac{y^2}{b^2}=1\) & \(\dfrac{xx_0}{a^2}-\dfrac{yy_0}{b^2}=1\) \\[6pt]
\(y^2=2px\) & \(yy_0=p(x+x_0)\) \\
\bottomrule
\end{longtable}
\endgroup
椭圆的\(c^2=a^2-b^2\)（\(a\ge b\)），双曲线的\(c^2=a^2+b^2\)，两者离心率均为\(e=c/a\)。上述抛物线的焦点为\((p/2,0)\)，准线为\(x=-p/2\)。

椭圆可表示为\(x=a\cos t,y=b\sin t\)；双曲线可表示为\(x=\pm a\cosh t,y=b\sinh t\)；抛物线可表示为\(x=pt^2/2,y=pt\)。椭圆面积为\(\pi ab\)，上述双曲线的渐近线为\(y=\pm(b/a)x\)。

若非退化二次曲线在齐次坐标中写为\(X^TMX=0\)、\(M=M^T\)，点\(P\)的极线为\(P^TMX=0\)；切点\(P\)在曲线上时，这就是该点的切线。

\newpage
\section{求导公式}

下表中的参数按相应函数的定义取值，\(a>0,a\ne1\)，实数幂的基本定义域取\(x>0\)。
\begingroup
\renewcommand{\arraystretch}{1.1}
\begin{longtable}{@{}p{.44\linewidth}p{.48\linewidth}@{}}
\toprule
\(f(x)\) & \(f'(x)\) \\
\midrule
\endhead
\(x^\alpha\) & \(\alpha x^{\alpha-1}\) \\
\(e^x,\ a^x\) & \(e^x,\ a^x\ln a\) \\
\(\ln x,\ \log_a x\) & \(1/x,\ 1/(x\ln a)\) \\
\(\sin x,\ \cos x\) & \(\cos x,\ -\sin x\) \\
\(\tan x,\ \cot x\) & \(1/\cos^2x,\ -1/\sin^2x\) \\
\(\arcsin x,\ \arccos x\) & \(1/\sqrt{1-x^2},\ -1/\sqrt{1-x^2}\) \\
\(\arctan x\) & \(1/(1+x^2)\) \\
\(\sinh x,\ \cosh x\) & \(\cosh x,\ \sinh x\) \\
\(\tanh x,\ \coth x\) & \(1/\cosh^2x,\ -1/\sinh^2x\) \\
\bottomrule
\end{longtable}
\endgroup
\begin{equation*}
 \begin{aligned}
 (uv)'&=u'v+uv',&
 \left(\frac uv\right)'&=\frac{u'v-uv'}{v^2},\\
 (g\circ f)'(x)&=g'(f(x))f'(x),&
 (f^{-1})'(y)&=\frac1{f'(x)}\quad(y=f(x),\ f'(x)\ne0).
 \end{aligned}
\end{equation*}
对参数曲线\(x=x(t),y=y(t)\)，若\(x'(t)\ne0\)，则\(dy/dx=y'(t)/x'(t)\)。

\section{积分公式}

以下\(C\)为任意常数，\(\alpha\ne-1\)、\(a>0,a\ne1\)，各式在表达式有意义的区间上使用。
\begin{equation*}
 \begin{aligned}
 \int x^\alpha\,dx&=\frac{x^{\alpha+1}}{\alpha+1}+C,&
 \int\frac{dx}{x}&=\ln|x|+C,\\
 \int e^x\,dx&=e^x+C,&
 \int a^x\,dx&=\frac{a^x}{\ln a}+C,\\
 \int\sin x\,dx&=-\cos x+C,&
 \int\cos x\,dx&=\sin x+C,\\
 \int\frac{dx}{1+x^2}&=\arctan x+C,&
 \int\frac{dx}{\sqrt{1-x^2}}&=\arcsin x+C.
 \end{aligned}
\end{equation*}
在相应微积分定理的条件下，
\begin{equation*}
 \begin{gathered}
 \int_a^b f(x)\,dx=F(b)-F(a)\quad(F'=f),\qquad
 \int u\,dv=uv-\int v\,du,\\
 \int f(\varphi(t))\varphi'(t)\,dt=F(\varphi(t))+C,\\
 \frac d{dx}\int_{u(x)}^{v(x)}f(t)\,dt
 =f(v(x))v'(x)-f(u(x))u'(x).
 \end{gathered}
\end{equation*}

\section{Taylor公式与几何累积}

在相应Taylor定理的条件下，
\begin{equation*}
 f(x)=\sum_{k=0}^n\frac{f^{(k)}(a)}{k!}(x-a)^k+R_n(x),
\end{equation*}
其中常用的两种余项形式为
\begin{equation*}
 R_n(x)=\frac{f^{(n+1)}(\xi)}{(n+1)!}(x-a)^{n+1},\qquad
 R_n(x)=\frac1{n!}\int_a^x(x-t)^n f^{(n+1)}(t)\,dt.
\end{equation*}
第一式中的\(\xi\)位于\(a,x\)之间。二阶局部模型为\(f(a+h)=f(a)+f'(a)h+f''(a)h^2/2+o(h^2)\)。

对上下边界为\(f\ge g\)的区域、连续截面积为\(A(x)\)的几何体，以及光滑曲线，常用累积公式为
\begin{equation*}
 \begin{aligned}
 S&=\int_a^b(f-g)\,dx,&
 V&=\int_a^b A(x)\,dx,\\
 L&=\int_\alpha^\beta\sqrt{x'(t)^2+y'(t)^2}\,dt,&
 S_{\mathrm{rot}}&=2\pi\int_\alpha^\beta r(t)\sqrt{x'(t)^2+y'(t)^2}\,dt.
 \end{aligned}
\end{equation*}
\(r(t)\)为到旋转轴的距离；曲线或曲面按恰好覆盖一次取参数范围。曲率在速度非零处为
\begin{equation*}
 \kappa=\frac{|x'y''-y'x''|}{(x'^2+y'^2)^{3/2}}.
\end{equation*}

\section{微分方程}

一阶线性初值问题\(y'+py=q\)、\(y(t_0)=y_0\)中，记\(P(t)=\int_{t_0}^tp(s)\,ds\)，则
\begin{equation*}
 y(t)=e^{-P(t)}\left(y_0+\int_{t_0}^t e^{P(s)}q(s)\,ds\right).
\end{equation*}
二阶常系数齐次方程\(y''+ay'+by=0\)的特征方程为\(r^2+ar+b=0\)，通解按根型为
\begin{equation*}
 \begin{cases}
 C_1e^{r_1t}+C_2e^{r_2t},&r_1\ne r_2\text{为实根},\\
 (C_1+C_2t)e^{rt},&r\text{为重根},\\
 e^{\alpha t}(C_1\cos\beta t+C_2\sin\beta t),&\alpha\pm i\beta\ (\beta>0).
 \end{cases}
\end{equation*}
非齐次方程在已知一个特解\(y_p\)后，其通解为\(y_p+y_h\)，其中\(y_h\)是齐次通解。
\endgroup

\chapter{基本计数公式}
\label{app:counting-formulas}

本附录集中整理计数公式，供学习与解题时查阅。每个公式都对应一种确定的计数对象：是否区分顺序、是否允许重复、对象与位置是否可辨别，应由问题中的实际安排决定。相应概念及推导将在第六章集中展开。

\section{加法原理与乘法原理}

若完成一项任务的所有方法分成\(k\)个互不重叠的类别，第\(i\)类有\(m_i\)种方法，则总方法数为
\begin{equation*}
 N=m_1+m_2+\cdots+m_k.
\end{equation*}
若一项任务需要依次完成\(k\)步，且对每一种已经完成的前序安排，第\(i\)步都有\(m_i\)种选择，则总方法数为
\begin{equation*}
 N=m_1m_2\cdots m_k.
\end{equation*}
分类中每种结果应当恰好计入一次；分步中每条完整的选择路径应当对应一种结果。若不同路径通向相同结果，就要进一步辨认它们被重复计算的次数。

两个有限集合的计数满足
\begin{equation*}
 |A\cup B|=|A|+|B|-|A\cap B|.
\end{equation*}
三个集合时，
\begin{equation*}
 \begin{aligned}
 |A\cup B\cup C|={}&|A|+|B|+|C|\\
 &-|A\cap B|-|A\cap C|-|B\cap C|+|A\cap B\cap C|.
 \end{aligned}
\end{equation*}
若总对象集合为\(\Omega\)，也可用\(|A|=|\Omega|-|\Omega\setminus A|\)从补集计数。

\section{排列与组合}

记\(n!=1\cdot2\cdots n\)，约定\(0!=1\)。本节\(n\)为正整数、\(r\)为非负整数，排列数和组合数公式中取\(0\le r\le n\)。

\begingroup
\renewcommand{\arraystretch}{1.2}
\begin{longtable}{@{}p{.64\linewidth}p{.30\linewidth}@{}}
\toprule
计数对象 & 方法数 \\
\midrule
\endhead
将\(n\)个不同对象排成一列 & \(n!\) \\
从\(n\)个不同对象中不重复地选\(r\)个，按顺序排列 & \(A_n^r=\dfrac{n!}{(n-r)!}\) \\[5pt]
从\(n\)个不同对象中选\(r\)个，不区分顺序 & \(C_n^r=\binom nr=\dfrac{n!}{r!(n-r)!}\) \\[5pt]
每个位置都可从\(n\)种对象中选一种，共\(r\)个有区别的位置，允许重复 & \(n^r\) \\
\bottomrule
\end{longtable}
\endgroup

排列与组合的关系为
\begin{equation*}
 A_n^r=C_n^r r!.
\end{equation*}
同一个\(r\)元子集具有\(r!\)种内部排列，所以无序选择与有序排列之间相差这一固定倍率。

若\(n\)个对象分成\(k\)类，同类对象不可辨别，第\(i\)类有\(n_i\)个，且\(n_1+\cdots+n_k=n\)，则不同的线性排列数为
\begin{equation*}
 \frac{n!}{n_1!n_2!\cdots n_k!}.
\end{equation*}
同一公式也计算将\(n\)个不同对象分配到\(k\)个有区别的组、各组人数依次为\(n_1,\ldots,n_k\)的方案数。

\section{二项式定理}

对非负整数\(n\)，
\begin{equation*}
 (a+b)^n=\sum_{k=0}^n\binom nk a^{n-k}b^k.
\end{equation*}
展开式中的第\(k+1\)项为\(\binom nk a^{n-k}b^k\)。若两个因子本身含有变量，所求变量次数应由整项中的幂次共同确定。

常用关系包括
\begin{equation*}
 \binom nk=\binom n{n-k},\qquad
 \binom nk=\binom{n-1}{k}+\binom{n-1}{k-1}.
\end{equation*}
求和时可约定\(\binom nk=0\)（\(k<0\)或\(k>n\)），使端部的公式统一。后一恒等式中取\(n\ge1\)。

代入不同的\(a,b\)，或对二项式展开式作有限次求导，可得
\begin{equation*}
 \begin{aligned}
 \sum_{k=0}^n\binom nk&=2^n,\\
 \sum_{k=0}^n(-1)^k\binom nk&=0\qquad(n\ge1),\\
 \sum_{k=0}^n k\binom nk&=n2^{n-1}\qquad(n\ge1),\\
 \sum_{k=0}^n k(k-1)\binom nk&=n(n-1)2^{n-2}\qquad(n\ge2).
 \end{aligned}
\end{equation*}
对非负整数\(m,n\)及\(0\le r\le m+n\)，还有
\begin{equation*}
 \sum_{k=0}^r\binom mk\binom n{r-k}=\binom{m+n}r.
\end{equation*}
它对应从两组共\(m+n\)个不同对象中选出\(r\)个，按从第一组选出的数量分类计数。

\section{分组与分配的补充公式}

\paragraph{圆排列} 将\(n\ge1\)个不同对象排成一圈，只把旋转视为相同、保留方向差别时，方法数为\((n-1)!\)。固定一个对象的位置后，其余对象按顺序排列即可。

\paragraph{重复选择} 从\(n\ge1\)类对象中选择\(r\ge0\)个，同类对象不区分、允许重复且不计顺序时，方法数为
\begin{equation*}
 \binom{n+r-1}r.
\end{equation*}
它等价于计算非负整数解\(x_1+\cdots+x_n=r\)的个数。

\paragraph{相同对象的分配} 将\(r\)个相同对象放入\(n\ge1\)个有区别的盒子，允许空盒时有\(\binom{r+n-1}{n-1}\)种；每盒至少一个且\(r\ge n\)时有\(\binom{r-1}{n-1}\)种。后一公式先给每盒放入一个，再分配剩余对象。

\paragraph{不同对象的分配} 将\(r\)个不同对象放入\(n\)个有区别的盒子，允许空盒时有\(n^r\)种。若规定每盒的对象数为\(r_1,\ldots,r_n\)，则有\(r!/(r_1!\cdots r_n!)\)种。

这些公式中的“不同”“相同”描述对象能否辨别，“有序”“无序”描述结果是否保留顺序，“盒子有区别”描述位置是否带有固定标记。它们决定同一结果应怎样识别，也决定是否需要除去重复计数的倍率。

\chapter{概率统计公式}
\label{app:probability-formulas}

本附录整理第六章拟使用的基本公式，供后续学习查阅；概念、推导和模型解释在相应正文中展开。计数公式另见附录\ref{app:counting-formulas}。

\section{事件与条件概率}
\begin{equation*}
 \begin{aligned}
 P(A^c)&=1-P(A),\\
 P(A\cup B)&=P(A)+P(B)-P(A\cap B),\\
 P(A\mid B)&=\frac{P(A\cap B)}{P(B)}\quad(P(B)>0),\\
 P(A\cap B)&=P(B)P(A\mid B).
 \end{aligned}
\end{equation*}
若\(B_1,\ldots,B_m\)为样本空间的有限分割，且各\(P(B_i)>0\)，则
\begin{equation*}
 P(A)=\sum_{i=1}^mP(B_i)P(A\mid B_i),\qquad
 P(B_j\mid A)=\frac{P(B_j)P(A\mid B_j)}{\sum_{i=1}^mP(B_i)P(A\mid B_i)}
\end{equation*}
（后一式取\(P(A)>0\)）。事件\(A,B\)独立的定义为\(P(A\cap B)=P(A)P(B)\)。

\section{期望与波动}

对离散变量\(P(X=x_i)=p_i\)，
\begin{equation*}
 EX=\sum_i x_ip_i,\qquad
 \operatorname{Var}X=E(X-EX)^2=E(X^2)-(EX)^2.
\end{equation*}
无穷取值情形的期望与方差按相应的存在条件使用。对具有密度\(f\)的模型，积分形式为\(Eg(X)=\int g(x)f(x)\,dx\)。
\begin{equation*}
 \begin{aligned}
 E(aX+bY)&=aEX+bEY,\\
 \operatorname{Var}(aX+b)&=a^2\operatorname{Var}X,\\
 \operatorname{Cov}(X,Y)&=E[(X-EX)(Y-EY)],\\
 \operatorname{Var}(X+Y)&=\operatorname{Var}X+\operatorname{Var}Y+2\operatorname{Cov}(X,Y),\\
 \rho_{XY}&=\frac{\operatorname{Cov}(X,Y)}{\sqrt{\operatorname{Var}X\operatorname{Var}Y}}.
 \end{aligned}
\end{equation*}
相关系数在两个方差都为正时使用；独立且二阶矩有限时，协方差为零。对非负变量及\(t>0\)，\(P(X\ge t)\le EX/t\)；对方差有限的变量及\(\varepsilon>0\)，
\begin{equation*}
 P(|X-EX|\ge\varepsilon)\le\frac{\operatorname{Var}X}{\varepsilon^2}.
\end{equation*}

\section{常用分布}

二项分布\(X\sim\mathrm B(n,p)\)表示\(n\)次独立Bernoulli试验的成功次数：
\begin{equation*}
 P(X=k)=\binom nk p^k(1-p)^{n-k},\quad 0\le k\le n;
 \qquad EX=np,\quad\operatorname{Var}X=np(1-p).
\end{equation*}
超几何模型中，总体有\(N\)个对象，其中\(M\)个具有指定性质；不放回抽取\(n\)个，令\(X\)记录其中具有该性质的数量，则
\begin{equation*}
 P(X=k)=\frac{\binom Mk\binom{N-M}{n-k}}{\binom Nn},\qquad EX=n\frac MN,
\end{equation*}
其中\(\max(0,n-N+M)\le k\le\min(n,M)\)。当\(N>1\)时，
\begin{equation*}
 \operatorname{Var}X=n\frac MN\left(1-\frac MN\right)\frac{N-n}{N-1}.
\end{equation*}
几何分布记录成功概率为\(0<p\le1\)的独立试验中首次成功的次数，取值从\(1\)开始：
\begin{equation*}
 P(X=k)=(1-p)^{k-1}p,\quad EX=\frac1p,\quad
 \operatorname{Var}X=\frac{1-p}{p^2}.
\end{equation*}
正态模型\(X\sim N(\mu,\sigma^2)\)、\(\sigma>0\)，密度与标准化为
\begin{equation*}
 f(x)=\frac1{\sigma\sqrt{2\pi}}e^{-(x-\mu)^2/(2\sigma^2)},\qquad
 Z=\frac{X-\mu}{\sigma}\sim N(0,1).
\end{equation*}
记标准正态分布函数为\(\Phi\)，则
\begin{equation*}
 P(a<X\le b)=\Phi\left(\frac{b-\mu}\sigma\right)-\Phi\left(\frac{a-\mu}\sigma\right).
\end{equation*}

\section{样本与回归}

对数据\(x_1,\ldots,x_n\)，描述均值和方差为
\begin{equation*}
 \bar x=\frac1n\sum_i x_i,\qquad
 s_n^2=\frac1n\sum_i(x_i-\bar x)^2=\frac1n\sum_i x_i^2-\bar x^2.
\end{equation*}
在独立同分布、均值为\(\mu\)、方差为\(\sigma^2\)的抽样模型下，
\begin{equation*}
 E\bar X=\mu,\qquad\operatorname{Var}\bar X=\frac{\sigma^2}{n},\qquad
 S^2=\frac1{n-1}\sum_i(X_i-\bar X)^2\quad(n\ge2)
\end{equation*}
中的\(S^2\)是总体方差的无偏估计量。它与用于描述已有数据的\(s_n^2\)分别按任务使用。

对成对数据\((x_i,y_i)\)，记
\begin{equation*}
 S_{xx}=\sum_i(x_i-\bar x)^2,\quad
 S_{yy}=\sum_i(y_i-\bar y)^2,\quad
 S_{xy}=\sum_i(x_i-\bar x)(y_i-\bar y).
\end{equation*}
样本相关系数与最小二乘直线\(\widehat y=\widehat a+\widehat b x\)为
\begin{equation*}
 r=\frac{S_{xy}}{\sqrt{S_{xx}S_{yy}}},\qquad
 \widehat b=\frac{S_{xy}}{S_{xx}},\qquad
 \widehat a=\bar y-\widehat b\bar x.
\end{equation*}
相关系数要求\(S_{xx},S_{yy}>0\)，回归斜率要求\(S_{xx}>0\)。拟合结果的统计解释还取决于数据的获取及所用模型。

\section{列联表统计量}

\(2\times2\)列联表的四格频数依次记为\(a,b,c,d\)，总数\(N=a+b+c+d\)。当各行、列合计均为正时，Pearson统计量为
\begin{equation*}
 \chi^2=\frac{N(ad-bc)^2}{(a+b)(c+d)(a+c)(b+d)}.
\end{equation*}
一般形式为\(\sum_{i,j}(O_{ij}-E_{ij})^2/E_{ij}\)，其中\(O_{ij}\)为观测频数，独立假设下的\(E_{ij}\)由行合计乘列合计再除以总数得到。检验时使用的近似分布、临界值及适用条件应结合相应抽样模型确定。

\clearpage
\endgroup
